% --- Άσκηση 38850 (38850.tex) ---
\Askhsh{38850}{4}
\begin{ylh}{2}
    \item 2.1. Οι Πράξεις και οι Ιδιότητές τους
\end{ylh}
Μια μικρή βιοτεχνία φτιάχνει κρίκους για πετσέτες, ανοίγοντας σε ξύλινες σφαίρες κυλινδρικές τρύπες με διάφορες διαμέτρους $\delta$ και ύψος $h = 4\ cm$ (με $\delta > h$). Στο παρακάτω σχήμα βλέπουμε δυο τέτοιους κρίκους. Αποτελείται κάποιος από τους παρακάτω κρίκους από περισσότερο ξύλο;

\begin{center}
\DrawSphereInCylinder{70}{2}{1}
\DrawSphereInCylinder{70}{2.5}{1.25}
\end{center}

Μπορείτε να ελέγξετε την απάντησή σας στην παραπάνω ερώτηση ως εξής:

Παρακάτω βλέπετε τους κρίκους που έχουν προκύψει από τις ξύλινες σφαίρες, από τις οποίες έχουν αφαιρεθεί δυο σφαιρικά τμήματα πάνω και κάτω (δυο καπέλα) και έχει ανοιχτεί μια κυλινδρική τρύπα. Πιο συγκεκριμένα, ο μικρός κρίκος (αριστερά) έχει προκύψει από ξύλινη σφαίρα ακτίνας $R = 2{,}5\ cm$ και ο μεγάλος δεξιά από ξύλινη σφαίρα ακτίνας $R' = 4\ cm$.

\begin{alist}
\item

\begin{rlist}
\item Να υπολογίσετε την ακτίνα της κυλινδρικής τρύπας σε κάθε ένα από τα παραπάνω σχήματα.
\monades{6}
\item Να υπολογίσετε τον όγκο της κυλινδρικής τρύπας σε κάθε περίπτωση και να τον γράψετε στη μορφή $V = a \cdot \pi$, με $a$ φυσικό αριθμό.
\monades{6}
\end{rlist}
\item

\begin{rlist}
\item Αν ο όγκος της κυλινδρικής τρύπας στον μικρό κρίκο είναι $V_A = 9\pi\ cm^{3}$ και στον μεγάλο κρίκο είναι $V_B = 48\pi\ cm^{3}$ και αν είναι γνωστό ότι ο όγκος των δύο σφαιρικών τμημάτων (τα δύο καπέλα πάνω και κάτω σε κάθε σχήμα) είναι για την μικρή σφαίρα $\dfrac{7}{6}\pi$ και για τη μεγάλη $\dfrac{80}{3}\pi$, να βρείτε τον όγκο του ξύλου που έμεινε (μετά τη δημιουργία της κυλινδρικής τρύπας) σε καθέναν από τους παραπάνω κρίκους.
\monades{8}
\item Στη βιοτεχνία οι τεχνίτες προβληματίστηκαν με το εξής: μήπως αν επιλέξουν μια οποιαδήποτε σφαίρα ακτίνας $R_1 > 2$, για να φτιάξουν τον κρίκο με τον ίδιο τρόπο που περιγράφεται παραπάνω, θα μείνει ο ίδιος όγκος ξύλου που έμεινε και στους δυο προηγούμενους; Ρώτησαν έναν φίλο τους που μπορούσε να τους βοηθήσει και τους είπε: «Ο όγκος του ξύλου που θα μείνει θα είναι:
\begin{equation*}
V = \dfrac{4}{3}\pi R_1^{3} - 4\pi\left(R_1^{2} - 4\right) - \dfrac{4}{3}\pi\left(R_1 - 2\right)^{2}\left(R_1 + 1\right)
\end{equation*}
και είναι ίσος με τον όγκο του ξύλου που έμεινε και στους δυο προηγούμενους (που βρήκατε στο βi) ερώτημα)».

Με δεδομένο τον παραπάνω τύπο, να ελέγξετε αν ο ισχυρισμός είναι σωστός. Μπορούμε να βγάλουμε ένα συμπέρασμα για τον όγκο του ξύλου που μένει σε κάθε περίπτωση στους κρίκους και τελικά να απαντήσουμε στην ερώτηση που τέθηκε στην αρχή;

\monades{5}

Δίνεται ο όγκος του κυλίνδρου: $V = \pi\rho^{2} \cdot h$, όπου $\rho$ η ακτίνα και $h$ το ύψος του κυλίνδρου και ο όγκος της σφαίρας $V = \dfrac{4}{3}\pi R^{3}$, όπου $R$ η ακτίνα της σφαίρας.
\end{rlist}
\end{alist}
